\documentclass[11pt,letterpaper]{article}
\usepackage[margin=1in]{geometry}       % 1-inch margins on all sides :contentReference[oaicite:1]{index=1}
\usepackage{graphicx}                   % Enhanced includegraphics interface :contentReference[oaicite:2]{index=2}
\usepackage{subcaption}

\usepackage{caption}                    % Customizable captions :contentReference[oaicite:5]{index=5}
\usepackage{float}                      % For [H] placement specifier
\usepackage{hyperref}                   % Clickable hyperlinks :contentReference[oaicite:6]{index=6}
\hypersetup{colorlinks=true,linkcolor=blue,urlcolor=blue}
\setlength{\parskip}{6pt}
\setlength{\parindent}{0pt}
\begin{document}

Overview:
\begin{itemize}
\item Similar results for Pangu and IFS
\item Might be able to improve forecasts more closer to the equator? Would like to run with larger patches...
\item Forecasts are generally better closer to the equator. This effect is larger for longer lead times. Is this since weather is more stable closer to the equator?
\item SDOR doesn't seem to matter for temp but can improve wind speed forecasts more in areas with high SDOR where baseline accuracy is lower. 
\end{itemize}

\section{Pangu: 2m Temperature}
\begin{figure}[H]
    \centering
    \includegraphics[width=\textwidth]{/Users/ohouck/globus/forecast_data/figures/pangu/global_maps/global_original_rmse_map_2m_temperature_pangu_lt216.png}
\end{figure}
\begin{figure}[H]
    \centering
    \includegraphics[width=\textwidth]{/Users/ohouck/globus/forecast_data/figures/pangu/global_maps/global_original_rmse_map_pixel_2m_temperature_pangu_lt216.png}
\end{figure}
\begin{figure}[H]
    \centering
    \includegraphics[width=\textwidth]{/Users/ohouck/globus/forecast_data/figures/pangu/global_maps/global_original_rmse_map_2m_temperature_pangu_lt24.png}
\end{figure}

\section{IFS: 2m Temperature}


\begin{figure}[H]
    \centering
    \includegraphics[width=\textwidth]{/Users/ohouck/globus/forecast_data/figures/ifs/global_maps/global_original_rmse_map_2m_temperature_ifs_lt216.png}
\end{figure}
\begin{figure}[H]
    \centering
    \includegraphics[width=\textwidth]{/Users/ohouck/globus/forecast_data/figures/ifs/global_maps/global_original_rmse_map_pixel_2m_temperature_ifs_lt216.png}
\end{figure}
\begin{figure}[H]
    \centering
    \includegraphics[width=\textwidth]{/Users/ohouck/globus/forecast_data/figures/ifs/global_maps/global_original_rmse_map_2m_temperature_ifs_lt24.png}
\end{figure}

\section{Pangu: 10m Wind Speed}
\begin{figure}[H]
    \centering
    \includegraphics[width=\textwidth]{/Users/ohouck/globus/forecast_data/figures/pangu/global_maps/global_original_rmse_map_10m_wind_speed_pangu_lt216.png}
\end{figure}
\begin{figure}[H]
    \centering
    \includegraphics[width=\textwidth]{/Users/ohouck/globus/forecast_data/figures/pangu/global_maps/global_original_rmse_map_pixel_10m_wind_speed_pangu_lt216.png}
\end{figure}
\begin{figure}[H]
    \centering
    \includegraphics[width=\textwidth]{/Users/ohouck/globus/forecast_data/figures/pangu/global_maps/global_original_rmse_map_10m_wind_speed_pangu_lt24.png}
\end{figure}


\section{IFS: 10m Wind Speed}

\begin{figure}[H]
    \centering
    \includegraphics[width=\textwidth]{/Users/ohouck/globus/forecast_data/figures/ifs/global_maps/global_original_rmse_map_10m_wind_speed_ifs_lt216.png}
\end{figure}
\begin{figure}[H]
    \centering
    \includegraphics[width=\textwidth]{/Users/ohouck/globus/forecast_data/figures/ifs/global_maps/global_original_rmse_map_pixel_10m_wind_speed_ifs_lt216.png}
\end{figure}
\begin{figure}[H]
    \centering
    \includegraphics[width=\textwidth]{/Users/ohouck/globus/forecast_data/figures/ifs/global_maps/global_original_rmse_map_10m_wind_speed_ifs_lt24.png}
\end{figure}


\end{document}